\subsection{Model Families and the Scheme of Elimination}
\label{sec:families}

Six model families are screened on the representation of Section~\ref{sec:representation}: the $k$-nearest-neighbour regressor and the local thin-plate interpolant of that section, Gaussian process regression~\citep{Rasmussen2005_GP} with a scaled squared-exponential kernel plus white noise, gradient-boosted trees (XGBoost~\citep{TODO_XGBoost}), the baseline multilayer perceptron, and a residual network of the same width and depth whose blocks add to identity skips. Each family runs in both units of prediction, one row at a time and curve-wise, where it learns the spline coefficients and the range of each curve from the curve parameters; the curve-wise interpolant is the interpolation between curves of the outline. The twelve candidates per dataset and target are eliminated by successive halving over the training rows: every candidate is scored through the harness of Section~\ref{sec:optimization} on $10^3$ training rows, the best three of each unit on the frozen folds go on to $10^4$, and the best two of each unit to $10^5$. Each unit is ranked apart: on $10^3$ rows the neutron-star design keeps about a hundred of its curves, too few for a curve-wise base, which at full data was the best representation of Section~\ref{sec:representation}. A smaller budget keeps rows evenly spaced along every curve, at least ten per curve, and fewer whole curves when the rows do not suffice for all. Each candidate has a fixed budget, its default settings and one seed; a fit past $30$ minutes or $4$\,GB is the family's wall at that number of rows, and a pointwise Gaussian process whose kernel matrix would pass the memory budget is not run. \Cref{tab:families} reports the significant figures each candidate keeps on the folds at the most rows it reached.

\InputIfFileExists{62_families_tab_families}{}{}

\todo{The results of the screen, from its table: the survivors per pair, set beside earlier neural surrogates~\citep{Liodis2024,Kaushikk2026,Liu2026}.}
\begin{itemize}
  \item $H_3$: the ripple of the residual across the grid steps of the curve parameters, a residual network against a plain multilayer perceptron of matched size.
  \item \textbf{Scalability:} the error and the fit time against the number of training rows, with the $O(N^3)$ wall of Gaussian process regression.
\end{itemize}
